\documentclass[11pt,a4paper]{article}

\usepackage{fontspec}
\usepackage{polyglossia}
\setdefaultlanguage{russian}
\setotherlanguage{english}

\setmainfont{Times New Roman}
\setsansfont{Arial}
\setmonofont{Courier New}

% Явно задаём кириллические варианты для polyglossia
\newfontfamily\cyrillicfont{Times New Roman}
\newfontfamily\cyrillicfontsf{Arial}
\newfontfamily\cyrillicfonttt{Courier New}

% ==================== ПАКЕТЫ ====================
\usepackage{geometry}
\geometry{margin=2.5cm}
\usepackage{enumitem}
\usepackage{booktabs}
\usepackage{hyperref}
\hypersetup{
    colorlinks=true,
    linkcolor=blue,
    urlcolor=blue,
}
\usepackage{xcolor}

% ==================== НАСТРОЙКИ ====================
\setlist[enumerate]{itemsep=4pt, topsep=4pt, leftmargin=1.5em}
\setlist[itemize]{itemsep=2pt, topsep=2pt, leftmargin=1.5em}
\setlength{\parindent}{0pt}
\setlength{\parskip}{6pt}

\newcommand{\code}[1]{\texttt{\colorbox{gray!10}{\detokenize{#1}}}}

\begin{document}

\begin{center}
    {\Large\bfseries Инструкция по созданию форка в GitHub}\\[4pt]
\end{center}

\vspace{0.5em}
\hrule
\vspace{0.5em}

\section*{Что нужно сделать}

\begin{enumerate}
    \item Открой ссылку на репозиторий:\\
          \url{https://github.com/vznuzdaev/Poster}
    \item Нажми кнопку \textbf{Fork} (справа сверху).\\
          Это создаст \textbf{твою личную копию} этого репозитория на GitHub.

    \item На странице своей копии нажми зелёную кнопку \textbf{Code} и скопируй ссылку.\\
          Затем на своём компьютере открой терминал и выполни:
          \begin{quote}
              \code{git clone <ссылка-на-твою-копию>}
          \end{quote}
          Это скачает копию проекта на твой компьютер.

    \item Зайди в папку с проектом:
          \begin{quote}
              \code{cd <название-папки>}
          \end{quote}

    \item Создай \textbf{отдельную ветку} для своих изменений (чтобы не трогать основную версию):
          \begin{quote}
              \code{git checkout -b fix/my-changes}
          \end{quote}

    \item Внеси изменения в файлы (редактором, как обычно).

    \item Сохрани изменения в истории проекта:
          \begin{quote}
              \code{git add .}\\
              \code{git commit -m "краткое описание, что изменил"}
          \end{quote}

    \item Отправь изменения в \textbf{свою копию} на GitHub:
          \begin{quote}
              \code{git push origin fix/my-changes}
          \end{quote}

    \item На странице своей копии на GitHub появится кнопка \textbf{Compare \& pull request}~--- нажми её.\\
          Убедись, что:
          \begin{itemize}
              \item слева (куда)~--- \textbf{мой} репозиторий, ветка \code{main};
              \item справа (откуда)~--- \textbf{твоя} ветка \code{fix/my-changes}.
          \end{itemize}
          Нажми \textbf{Create pull request}.

    \item Я получу уведомление, посмотрю изменения и, если всё в порядке, добавлю их в основной репозиторий.
\end{enumerate}

\end{document}